% !TEX root = template.tex

\section{Concluding Remarks}
\label{sec:conclusions}

We evaluated, under one shared apparatus and a single-use test protocol,
whether adversarial (GRL) and contrastive (SupCon) objectives improve
cross-environment WiFi CSI activity recognition over plain cross-entropy. They
do not: SupCon is worse by a bootstrap-resolved margin and the GRL is worse
still and less stable across seeds, while the deep CE model beats the
shallow--wide reference design by $23.5$ accuracy points --- backbone capacity
mattered far more than any invariance objective. The observation we consider most useful is \emph{why} the
objectives fail: frozen-encoder probes show that no domain attribute is
linearly readable from the learned features, so a domain adversary has nothing
to remove --- a precondition that costs one linear probe on cached features to
check, and is worth checking on any dataset whose training split, like this
one, is dominated by a single environment --- a message that goes to benchmark
designers as much as to practitioners.

\subsection{Challenges and Lessons Learned}

Much of this project lay outside the material we had seen: test-time
adaptation, the trace-level paired bootstrap and the frozen-encoder probing
protocol were designed from the literature rather than assembled from known
recipes, and while contrastive learning was familiar in its self-supervised
form, its supervised variant and the gradient reversal layer had to be
implemented and debugged from the original papers. The hardest decisions came
from the data, imbalanced along every axis at once --- environments, capture
sets, subjects and classes --- so that each remedy costs something: pinning
rare cells into training guarantees coverage but empties part of validation,
balancing contrastive batches over capture sets oversamples trace-poor sets,
and a third rotation that looked mandatory on paper had to be rejected for
having $26$ training traces. We also learned to distrust single numbers: the
seed replicates showed a same-configuration swing on the GRL larger than some
cross-configuration gaps, and with $11$ test traces even a $16$-point gap is
unresolved once the honest unit of analysis is the trace. Finally, a caveat we
had declared in advance --- that classes invisible to validation would suffer
--- proved false on test, a reminder that plausible methodological worries
should be measured, not just declared.

\textbf{Future work.} With one held-out session per rotation this is a
transfer study rather than full domain generalization; the natural extension
is a dataset with several well-populated training environments, where our
diagnostic protocol transfers unchanged. Whether \emph{any} training-time
objective can help when the domain factor is unreadable remains open ---
candidates acting on the input, such as physics-informed augmentation, are the
ones unaffected by our structural argument.
